\chapitresansnumero{Introduction générale}
Une association étudiante produit un journal numérique avec plusieurs centaines de rédacteurs. Ils doivent coéditer, consulter l'historique et restaurer une version. Des liens anonymes donnent accès à un nombre limité de PDF.


Le système doit garantir le quota, conserver chaque modification acceptée, tenir le crédit quotidien et produire un PDF cohérent malgré la concurrence.


Ce rapport présente les besoins, les modèles UML et les choix de conception. Les résultats expérimentaux relèvent du rapport de test.

Il comprend trois parties : état de l'art, analyse des besoins et conception.
